\documentclass{article}
\usepackage{amsmath}
\usepackage{amssymb}
\usepackage{geometry}
\geometry{a4paper, margin=1in}

\begin{document}

\title{Diffusion and Scaling in the Non-Linear Committee Machine}
\author{}
\date{}
\maketitle

\section{Introduction}
We analyze the scaling behavior and weight space dynamics of the non-linear committee machine learning a teacher generated by the third Hermite polynomial, $y = \text{He}_3(\boldsymbol{w}_* \cdot \boldsymbol{x})$. We investigate the typical scales of the overlaps, the resulting gradients, and the diffusion motion in the weight space during learning (e.g., via Langevin dynamics or gradient descent). 

\section{Scales of the Expressions}
At initialization, the student weights are drawn from an isotropic Gaussian prior:
\[
\boldsymbol{w}_i \sim \mathcal{N}\left(\boldsymbol{0}, \frac{1}{d}\boldsymbol{I}_d\right).
\]
In the high-dimensional limit ($d \to \infty$), the norm of the weights concentrates:
\[
\|\boldsymbol{w}_i\|^2 \approx 1.
\]
The teacher's weight vector is fixed and normalized, $\|\boldsymbol{w}_*\|^2 = 1$. The overlap between the student's hidden unit $i$ and the teacher is defined as:
\[
m_i = \boldsymbol{w}_i \cdot \boldsymbol{w}_*.
\]
Since $\boldsymbol{w}_i$ is a random vector with variance $1/d$ in each direction, the initial overlap is a sum of independent random variables:
\[
m_i = \sum_{k=1}^d w_{ik} w_{*k}.
\]
The variance of this initial overlap is $\frac{1}{d} \|\boldsymbol{w}_*\|^2 = \frac{1}{d}$. Therefore, the typical scale of the overlap at initialization is governed by random fluctuations:
\[
m_i \sim \mathcal{O}\left(\frac{1}{\sqrt{d}}\right).
\]

\section{The Signal and the $1/d$ Scaling}
You proposed that the scale of $\boldsymbol{w} \cdot \boldsymbol{w}_*$ goes like $1/d$. This theory is exactly correct when we consider the \textbf{gradient signal} or \textbf{drift} of the overlap, rather than its random initialization magnitude. Let us see why.

The student-teacher overlap term in the Hamiltonian (or expected loss) is $M_a = \mathbb{E}_{\boldsymbol{x}}[f(\boldsymbol{x}; \boldsymbol{W}) y(\boldsymbol{x})]$. For the error-function committee and the $\text{He}_3$ teacher, this evaluates to:
\[
M_a \propto \sum_{i=1}^N \left( \boldsymbol{w}_i \cdot \boldsymbol{w}_* \right)^3 = \sum_{i=1}^N m_i^3.
\]
During gradient descent or Langevin dynamics, the weights evolve according to the negative gradient of the loss. The gradient of the overlap term $M_a$ with respect to $\boldsymbol{w}_i$ provides the ``signal'' that pushes the student towards the teacher:
\[
\nabla_{\boldsymbol{w}_i} M_a \propto 3 m_i^2 \boldsymbol{w}_*.
\]
To see how this affects the overlap $m_i$ in one step of learning, we project the gradient update onto the teacher direction $\boldsymbol{w}_*$:
\[
\Delta m_i = (\Delta \boldsymbol{w}_i) \cdot \boldsymbol{w}_* \propto (\nabla_{\boldsymbol{w}_i} M_a) \cdot \boldsymbol{w}_* \propto m_i^2 \|\boldsymbol{w}_*\|^2 = m_i^2.
\]
Since the initial overlap is $m_i \sim \mathcal{O}(1/\sqrt{d})$, the gradient signal driving the overlap is:
\[
\Delta m_i \propto \left( \frac{1}{\sqrt{d}} \right)^2 = \mathcal{O}\left(\frac{1}{d}\right).
\]
Thus, your theory is verified: while the random, non-directional overlap is $1/\sqrt{d}$, the \textbf{directional signal (drift) pushing the weights toward the teacher scales exactly as $1/d$}. Because this signal is so weak in high dimensions (being overwhelmed by the $1/\sqrt{d}$ noise), the network experiences a long ``plateau'' or search phase before it can discover the teacher's direction.

\section{Diffusion Motion in the Weights Space}
In Langevin dynamics, the update rule for the weights is:
\[
d\boldsymbol{w}_i = -\nabla_{\boldsymbol{w}_i} H(\boldsymbol{W}) dt + \sqrt{2T} d\boldsymbol{B}_i,
\]
where $d\boldsymbol{B}_i$ is a standard Brownian motion in $\mathbb{R}^d$, and $T$ is the temperature.

If we look at the evolution of the weights orthogonal to the teacher, the dynamics are dominated by the Brownian noise and the regularization (or diagonal terms of the loss Hessian), leading to a standard diffusion process in the $d$-dimensional space. 

The equation for the overlap $m_i = \boldsymbol{w}_i \cdot \boldsymbol{w}_*$ becomes a scalar stochastic differential equation:
\[
dm_i = \left( - \lambda m_i + c m_i^2 + \dots \right) dt + \sqrt{\frac{2T}{d}} dB_t,
\]
where $\lambda$ comes from the weight decay and the student-student interaction $Q_{ab}$, and the signal $c m_i^2$ comes from the teacher overlap $M_a$.

Because the noise term $\sqrt{2T/d}$ is of order $1/\sqrt{d}$ and the drift signal is of order $1/d$, the initial phase of learning is heavily dominated by \textbf{diffusion}. The weights wander randomly in the high-dimensional space until a large enough fluctuation pushes $m_i$ to an $\mathcal{O}(1)$ value, at which point the $m_i^2$ drift term takes over and the student rapidly aligns with the teacher (escaping the plateau).

\section{Scaling of the Loss}
We can also calculate the scaling of the loss (the energy) at initialization. The energy is given by:
\[
E = \frac{1}{\kappa} \sum_{\mu=1}^P \left( f(\boldsymbol{x}^\mu; \boldsymbol{W}) - y^\mu \right)^2.
\]
Replacing the empirical average with the population expectation, the loss is proportional to the diagonal of the residual covariance matrix:
\[
\mathbb{E}[E] \approx \frac{\alpha d}{\kappa} \left( Q_{aa} - 2M_a + \rho \right).
\]
Let us analyze the scale of each term at initialization:
\begin{enumerate}
    \item \textbf{Teacher variance $\rho$}: For $\text{He}_3$, $\rho = \mathbb{E}[\text{He}_3(z)^2] = 6$. This is an $\mathcal{O}(1)$ constant.
    \item \textbf{Student self-overlap $Q_{aa}$}: As derived previously, $Q_{aa} \approx \frac{2}{\pi} \arcsin(2/3) \approx 0.46$. This is also an $\mathcal{O}(1)$ constant.
    \item \textbf{Student-teacher overlap $M_a$}: $M_a \propto \sum_{i=1}^N m_i^3$. Since $m_i \sim \mathcal{O}(1/\sqrt{d})$, $M_a \sim N \mathcal{O}(d^{-3/2})$. 
\end{enumerate}
Therefore, at the beginning of training, the loss is dominated by $Q_{aa} + \rho$, and the teacher signal $M_a$ provides only an $\mathcal{O}(d^{-3/2})$ perturbation to the macroscopic loss. The total loss scales as $\mathcal{O}(d)$ due to the extensive sum over $P = \alpha d$ examples, but the \emph{per-example} loss is $\mathcal{O}(1)$, with the information about the teacher buried in $1/d^{1.5}$ fluctuations.

\end{document}